\documentclass[11pt]{article}
\usepackage[a4paper,margin=27mm]{geometry}
\usepackage{amsmath,amssymb,amsthm}
\usepackage[hidelinks]{hyperref}
\newtheorem{theorem}{Theorem}
\newtheorem{lemma}{Lemma}
\newcommand{\Ie}{I_{\mathrm e}}
\title{\large Connected graphs with edge-irregulator number\texorpdfstring{ $3|E|/8$}{ 3|E|/8}}
\author{Research manuscript}
\date{29 October 2026 --- candidate; final review pending}
\begin{document}
\maketitle
\begin{abstract}
We construct, for every integer $t\ge2$, a connected simple graph with
$6t+1$ vertices, $8t$ edges and edge-irregulator number exactly $3t$.
The excess above one third of its edge count is therefore $t/3$.
This refutes Conjecture~1 of Bensmail, Catherinot, Fioravantes, Marcille
and Oijid, which asks for a uniform additive constant on connected graphs.
A two-part local charging argument gives the lower bound; an explicit
three-edge deletion in each block attains it. The proof is independent of
computer enumeration. A source-only Lean certificate accompanies this draft.
\end{abstract}

\section{Definition, source question and result}
All graphs below are finite, undirected and simple. A graph $H$ is
\emph{locally irregular} if $d_H(u)\ne d_H(v)$ for every edge $uv\in E(H)$.
For $S\subseteq E(G)$, let $G-S$ retain all vertices and remove the edges
in $S$. An \emph{edge-irregulator} is such an $S$ for which $G-S$ is locally
irregular, and
\[
 \Ie(G)=\min\{|S|:S\subseteq E(G),\ G-S\text{ is locally irregular}\}.
\]
The minimum exists since deleting all edges is allowed. In particular, the
endpoint degrees in this definition are \emph{recomputed in $G-S$}.
Neither connectivity nor positive minimum degree of $G-S$ is required.

Conjecture~1 of \cite[p.~356]{source} states that there exists an absolute
real $c\ge1$ such that every connected graph $G$ satisfies
\begin{equation}\label{conj}
 \Ie(G)\le |E(G)|/3+c.
\end{equation}
This is the publisher's connected-graph claim, with no assumption on the
original minimum degree. It is distinct from assertions restricted to
regular graphs or other special classes.

\begin{theorem}\label{main}
For every integer $t\ge2$ there is a connected simple graph $G_t$ such that
\[
 |V(G_t)|=6t+1,\qquad |E(G_t)|=8t,\qquad \Ie(G_t)=3t.
\]
Consequently no real constant $c\ge1$ satisfies \eqref{conj} for every
connected finite simple graph.
\end{theorem}

\newpage
\section{Construction and the lower bound}
Take one vertex $h$. For each $i\in\{1,\ldots,t\}$ take six private vertices
$a_i,b_i,c_i,d_i,x_i,y_i$. All these vertices are distinct, including across
blocks. The edges are precisely the following eight edges for each $i$:
\begin{equation}\label{edges}
 a_i b_i,\ a_i c_i,\ a_i d_i,\ b_i c_i,\ b_i d_i,
 \ c_i x_i,\ x_i y_i,\ y_i h.
\end{equation}
The first five edges form a diamond ($K_4$ minus $c_id_i$); the last three
form a path from $c_i$ to $h$. Every vertex has a path to $h$, so $G_t$ is
connected. There are no loops or duplicate edges. Different blocks share
only $h$, and their edge sets are disjoint. This proves the two claimed
cardinalities.

Fix an arbitrary edge-irregulator $S$, and put $H=G_t-S$. Fix a block and
omit its subscript. We make two disjoint charges against $S$.

\begin{lemma}\label{diamond}
At least two of $ab,ac,ad,bc,bd$ belong to $S$.
\end{lemma}
\begin{proof}
The only edge leaving the diamond is $cx$. Let $\varepsilon=1$ if $cx$
is retained and $\varepsilon=0$ otherwise. The degrees of $a,b,d$ in $H$
are exactly their retained diamond degrees, and the degree of $c$ is its
retained diamond degree plus $\varepsilon$.

If no diamond edge is deleted, $d_H(a)=d_H(b)=3$ contradicts the retained
edge $ab$. If exactly one is deleted, there are five possibilities.
The swap $a\leftrightarrow b$ preserves the diamond and its attachment at
$c$, so the following three cases include all five possibilities:
\begin{enumerate}
\item The deleted edge is $ab$. Then $a,b,d$ each have degree $2$.
The retained edge $ad$ has equal endpoint degrees.
\item The deleted edge is $ad$ (or, by the swap, $bd$).
Then $d_H(a)=2$, $d_H(b)=3$, and $d_H(c)=2+\varepsilon$.
For $\varepsilon=0$, retained $ac$ has degree pair $(2,2)$;
for $\varepsilon=1$, retained $bc$ has degree pair $(3,3)$.
\item The deleted edge is $ac$ (or, by the swap, $bc$).
Then $d_H(a)=d_H(d)=2$, contradicting the retained edge $ad$.
\end{enumerate}
Each case violates local irregularity, for both possible attachment states.
These degrees are those of the whole graph $H$; no exterior deletion can
alter them except through the already treated edge $cx$.
\end{proof}

\begin{lemma}\label{path}
At least one of $cx,xy,yh$ belongs to $S$.
\end{lemma}
\begin{proof}
Otherwise $x$ and $y$ each have degree $2$ in $H$, contradicting the retained
edge $xy$. The final degree of $h$ is irrelevant.
\end{proof}

\newpage
\section{Sharpness and the source quantifiers}
The two sets charged in Lemmas~\ref{diamond} and~\ref{path} are disjoint.
Moreover, all eight-edge block sets in \eqref{edges} are pairwise disjoint
and partition $E(G_t)$. Thus the charges count actual deleted edges without
multiplicity. For every admissible $S$,
\[
 |S|=\sum_{i=1}^t |S\cap E_i|\ \ge\ \sum_{i=1}^t(2+1)=3t,
\]
where $E_i$ is the eight-edge set of block $i$. This proves
$\Ie(G_t)\ge3t$. The sharing of $h$ does not identify any charged edges.
This argument works against arbitrary deletions, including ones that
isolate vertices or disconnect the retained graph.

For the reverse inequality, delete precisely
\[
 S_t=\bigcup_{i=1}^t\{a_i b_i,a_i d_i,x_i y_i\}.
\]
The sets in this union are disjoint, so $|S_t|=3t$. The private degrees in
$G_t-S_t$, in the order $(a_i,b_i,c_i,d_i,x_i,y_i)$, are
\[
 (1,2,3,1,1,1),\qquad d_{G_t-S_t}(h)=t.
\]
Every retained edge is one of the following five types; the table gives
its two final degrees:
\begin{center}
\begin{tabular}{c|ccccc}
edge & $a_ic_i$ & $b_ic_i$ & $b_id_i$ & $c_ix_i$ & $y_ih$\\\hline
final degrees & $(1,3)$ & $(2,3)$ & $(2,1)$ & $(3,1)$ & $(1,t)$
\end{tabular}
\end{center}
Every pair has unequal entries for $t\ge2$. Therefore $S_t$ is an
edge-irregulator and $\Ie(G_t)\le3t$, establishing equality.
The restriction $t\ge2$ is used in the last pair. No optimum for $t=1$
is asserted here.

To check the full negation of the source claim, let an arbitrary real
$c\ge1$ be proposed. By the Archimedean property choose a natural number
$t>3c$. Then $t\ge2$, and
\[
 \Ie(G_t)-\frac{|E(G_t)|}{3}
   =3t-\frac{8t}{3}=\frac{t}{3}>c.
\]
This contradicts \eqref{conj}. Thus the construction defeats every proposed
absolute additive constant, rather than only a particular numerical value.
Equivalently the excess is unbounded along this connected family.
\qed

\paragraph{Why the boundary calculation matters.}
An obstruction inside an induced subgraph need not remain an obstruction
after edges to its exterior change the endpoint degrees. The diamond
argument explicitly accommodates the only exterior contribution
$\varepsilon$, while the path argument avoids the hub degree completely.
This independence is what permits summing the lower bound over arbitrarily
many blocks joined through one hub. The finite cases above supply the
ordinary proof; an enumeration of masks is not a premise.

\newpage
\section{Formal certificate and status}
\raggedright
The accompanying certificate uses Lean 4.34.1 and pinned Mathlib with eight
other pinned packages. The thirteen inherited mathematical modules are
preserved byte for byte from the completed lower/source and upper-witness
components. The added \texttt{Result.lean} combines them into the closed
endpoints \texttt{EdgeIrregulator.minimum\_exact} and
\texttt{EdgeIrregulator.main}. The exact equality has only the hypotheses
$t:\mathbb N$ and $2\le t$.

\texttt{Source} defines actual graph edge deletion, the retained-degree
predicate and the attained minimum. \texttt{Gadget}, \texttt{Edges} and
\texttt{Degrees} define the graph and its edge indexing and neighbors.
\texttt{Boundary}, \texttt{Local}, \texttt{Lower} and \texttt{Charges}
connect the local constraints to arbitrary valid deletion sets and to
pairwise disjoint actual edge charges. The local finite certificate is
kernel checked. The self-contained cases on page 2 explain its mathematics.
\texttt{UpperWitness}, \texttt{UpperDegrees} and \texttt{Upper} verify
$S_t$, its cardinality and every retained degree and edge type.

\texttt{Transport} proves \texttt{minimum\_iso} for finite graph
isomorphisms, and transports $G_t$ to a graph on a finite ordinal.
The finite-model equality is the declaration
\texttt{EdgeIrregulator.fin\_minimum\_exact}; its module is \texttt{Result}.
\texttt{Refutation} proves \texttt{source\_refutation}, the negation of the
literal claim with one real $c\ge1$, every natural $n$, every connected
simple graph on $\mathrm{Fin}\,n$, and existence of a deletion set with
the claimed real bound. The closed main theorem includes connectivity,
vertex and edge counts, both exact minima and this source negation.

The exhaustive audit enumerates all 167 declarations from fourteen
mathematical module origins, including private and generated names. It
records full types, transitive axiom dependencies and safety flags, and
crosschecks the module declaration tables. Only \texttt{propext},
\texttt{Classical.choice} and \texttt{Quot.sound} are permitted.
Audit metaprograms are outside the mathematical module set. The packet
verifier checks the binary, all nine package revisions and imported
package hashes before a serial strict rebuild from source.

\paragraph{Attribution and limits.}
The author line follows the campaign's established ``Research manuscript''
attribution; no personal identity is inferred. OpenAI Codex was used for
proof and manuscript preparation, implementation and verification.
This is a candidate draft, not an accepted or published paper. Candidate
credit remains $0$; campaign counts remain $7/30$ additional and $37/60$
total. One family is proposed, with no separate credit for helpers or
individual parameters. Independent final root review and its final current
prior-art refresh remain pending.

The existing bounded source/prior review dated 29 October 2026 examined
the publisher definition and Conjecture~1, current registry inventories
and selected actual claims. It found no same-object counterpart or complete
subsumption bridge in that scope. It did not read every registry body or
certify worldwide priority; unindexed and unpublished work remain outside
its coverage. Exact review records, hashes, historical failures and resource
limitations are included with the certificate. No new novelty clearance
is claimed by this manuscript preparation.

\begin{thebibliography}{1}
\bibitem{source} J. Bensmail, N. Catherinot, F. Fioravantes, C. Marcille
and N. Oijid, \emph{Graph Irregularity via Edge Deletions}, WALCOM 2026,
pp.~353--367. Conjecture~1, p.~356.
\url{https://doi.org/10.1007/978-981-95-7127-7_24}.
Preprint: arXiv:2511.14514v1 (18 November 2025).
\end{thebibliography}
\end{document}
